\documentclass[12pt]{article}

% --- Encoding / language (XeLaTeX-safe) ---
\usepackage{fontspec}
\usepackage[margin=1in]{geometry}
\usepackage{amsmath, amsfonts, amssymb}
\usepackage{amsthm}
\numberwithin{equation}{section}
\usepackage[backend=biber, style=numeric, sorting=none]{biblatex}
\addbibresource{references.bib}

\newtheorem{proposition}{Proposition}
\theoremstyle{remark}
\newtheorem{remark}{Remark}

\begin{document}

\title{Kramers--Smoluchowski Limit: From Underdamped to Overdamped Langevin Dynamics}
\author{J.-F. Jabir \textit{for} D. Polozkov}
\date{October 9, 2024}
\maketitle

\section{Framework: Langevin dynamics}

Langevin dynamics define a class of second-order SDEs characterized by the following (formal) stochastic Newton law: 
\[ 
\gamma \ddot{X}_t^\gamma = F(t, X_t^\gamma, \dot{X}_t^\gamma)\,dt + \sigma(t, X_t^\gamma, \dot{X}_t^\gamma)\,dL_t, \qquad t \ge 0. 
\] 
From a physical point of view, $X^\gamma = (X_t^\gamma)_{t\ge0}$ models the position of a generic body evolving in $\mathbb{R}^d$, $\dot{X}^\gamma = (\dot{X}_t^\gamma)_{t\ge0}$ its velocity, $\gamma$ the volumic mass of the body, $F: [0,\infty)\times \mathbb{R}^d \times \mathbb{R}^d \to \mathbb{R}^d$ an external force field and $(L_t)_{t\ge0}$ an $m$-dimensional stochastic perturbation with dispersion matrix $\sigma$. Setting $V^\gamma = \dot{X}^\gamma$, the above rewrites more rigorously as the $\mathbb{R}^{2d}$-valued SDE system: 
\begin{equation}\label{eq:Lang}
\begin{aligned}
dX_t^\gamma &= V_t^\gamma\,dt, \\
dV_t^\gamma &= \gamma^{-1} F(t, X_t^\gamma, V_t^\gamma)\,dt + \gamma^{-1}\sigma(t, X_t^\gamma, V_t^\gamma)\,dL_t, \qquad t \ge 0~.
\end{aligned}
\end{equation} 
In this formulation the coefficient $\gamma^{-1}$ is understood as a friction force.

\begin{remark}
A prototypical instance of \eqref{eq:Lang} is given by the following setting: 
\begin{equation}\label{eq:protoSetting}
\begin{aligned}
L &= W \text{ (a $d$-dimensional Brownian motion)},\\
F(t,x,v) &= -\nabla f(x)\;-\;\lambda\,v,\\
\sigma(t,x,v) &= \sigma\,I_{d\times d} > 0~,
\end{aligned}
\end{equation}
with $\sigma$ a positive constant. In this situation, \eqref{eq:Lang} becomes 
\begin{equation}\label{eq:OU}
\begin{aligned}
dX_t^\gamma &= V_t^\gamma\,dt, \\
dV_t^\gamma &= -\,\gamma^{-1}\!\big(\nabla f(X_t^\gamma) + \lambda\,V_t^\gamma\big)dt + \gamma^{-1}\sigma\,dW_t, \qquad t \ge 0~,
\end{aligned}
\end{equation}
which is a Langevin equation often referred to as the underdamped Ornstein–Uhlenbeck process.
\end{remark}

Langevin dynamics, originally introduced in \cite{l08}, are applied in modern science across a wide range of fields, including fluid dynamics, biology, economics, and social science \cite{npt10,pt13}, and more recently in artificial intelligence \cite{ccb18}.

\section{Clarification on Pavliotis's scaling and limits}

(This section mainly refers to Chapter 6 in \cite{p14}.) The Langevin model in \cite{p14} (see Equation (6.1), p. 181) corresponds to a second-order stochastic system describing the motion of a “particle” with 
\[ 
dq_t = -\,\nabla U(q_t)\,dt \;-\; \eta\,\dot{q}_t\,dt \;+\; \sqrt{\frac{2\eta}{\beta}}\,dW_t, \qquad t \ge 0, 
\] 
where $q = (q_t)_{t\ge0}$ denotes the position process of the particle and $\dot{q} = (\dot{q}_t)_{t\ge0}$ the associated velocity process. Motions are assumed to start at the position $q_0$ and velocity $\dot{q}_0$. (For notational consistency, the potential force and friction coefficient are denoted here by $\nabla U$ and $\eta$ instead of $\nabla V$ and $\gamma$ as in \cite{p14}.) To this model, we associate the family of rescaled dynamics 
\[ 
q_t^\eta := \frac{\lambda_\eta}{\mu_\eta}\,q_{t/\mu_\eta}, \qquad t \ge 0. 
\] 
As such, the corresponding velocity is $\dot{q}_t^\eta = \frac{\lambda_\eta}{\mu_\eta}\,\dot{q}_{t/\mu_\eta}$, and so 
\begin{align*}
\dot{q}_t^\eta &= \frac{\lambda_\eta}{\mu_\eta}\,\dot{q}_0 \;-\; \frac{\eta\,\lambda_\eta}{\mu_\eta}\int_{0}^{\,t/\mu_\eta} \nabla U(q_s)\,ds \;-\; \frac{\lambda_\eta}{\mu_\eta}\int_{0}^{\,t/\mu_\eta} \dot{q}_s\,ds \;+\; \frac{\lambda_\eta}{\mu_\eta}\sqrt{\frac{2\eta}{\beta}}\,W_{t/\mu_\eta} \\[6pt]
&= \frac{\lambda_\eta}{\mu_\eta}\,\dot{q}_0 \;-\; \frac{\lambda_\eta}{\mu_\eta^2}\int_{0}^{\,t} \nabla U\!\big(q_{s/\mu_\eta}\big)ds \;-\; \frac{\eta\,\lambda_\eta}{\mu_\eta^2}\int_{0}^{\,t} \dot{q}_{s/\mu_\eta}\,ds \;+\; \frac{\lambda_\eta}{\mu_\eta}\sqrt{\frac{2\eta}{\beta}}\,W_{t/\mu_\eta}\,. 
\end{align*}
In the last step above, we changed variables in the integrals (set $u = s/\mu_\eta$) so that, for example, $\int_0^{\,t/\mu_\eta}\nabla U(q_s)ds = \mu_\eta^{-1}\int_0^t \nabla U(q_{u/\mu_\eta})du$. This formulation shows how time and space are rescaled in the underdamped-to-overdamped limit.

\section{Quantitative Kramers--Smoluchowski limit}

In the proposition below, under the prototypical condition \eqref{eq:protoSetting} (i.e. the setting of Remark 1), we establish the so-called \emph{Kramers--Smoluchowski limit}, which asserts the $\gamma \downarrow 0$ (zero-mass) limit of the system \eqref{eq:Lang}.

\begin{proposition}
Let $(X_t^\gamma, V_t^\gamma)$ be the Langevin dynamics given by \eqref{eq:OU}. Assume that the initial conditions $(X_0^\gamma, V_0^\gamma)$ are independent of $\gamma$ and equal to $(X^0, V^0)$, with $\mathbb{E}[|X^0|^2 + |V^0|^2] < \infty$, and that $\sigma \neq 0$ and $\nabla f$ is Lipschitz continuous. Then, for any finite time horizon $T$, $\gamma \le 1$ and any integer $p \ge 1$, 
\begin{equation}\label{eq:3.4}
\max_{\,t \in [0,T]}\;\mathbb{E}\!\Big[\,\big|X_t^\gamma - X_t^0\big|^{2p}\Big] = O(\gamma^p)\,.
\end{equation}
where $X^0 = (X_t^0)_{t\ge0}$ is the solution of the limiting SDE 
\begin{equation}\label{eq:3.5}
dX_t^0 = -\,\nabla f(X_t^0)\,dt + \sigma\,dW_t, \qquad X_0^0 = X^0~. 
\end{equation}
Moreover, 
\begin{equation}\label{eq:3.6}
\mathbb{E}\!\Big[\,\max_{\,t \in [0,T]}\;\big|X_t^\gamma - X_t^0\big|^2\Big] = O(\gamma^{1/2})\,. 
\end{equation}
\end{proposition}

While the Kramers--Smoluchowski limit is more directly linked to \eqref{eq:3.6}, this proposition is purposely written to emphasize the successive “easy” part \eqref{eq:3.4}\footnote{Stated without proof details in \cite{p14}, p. 206.} and the more “tricky” part \eqref{eq:3.6} of the limit. 

\begin{remark}
For $p = 1$, the estimate \eqref{eq:3.4} is optimal, but the discrepancy between the rate $O(\gamma)$ in \eqref{eq:3.4} and the $O(\gamma^{1/2})$ in \eqref{eq:3.6} is rather non-intuitive (and possibly subject to refinement). This discrepancy reflects the cost difference between time-marginal and max-in-time mean-square distances (which may be significant).
\end{remark}

\begin{proof}
As a preliminary, recall that under the initial and Lipschitz conditions of the proposition, each of SDEs \eqref{eq:Lang} and \eqref{eq:3.5} admits a unique strong solution with finite second moments at all times. In particular, since $\nabla f$ is Lipschitz, it grows at most linearly: 
\begin{equation}\label{eq:lipBound}
|\nabla f(x)| \;\le\; |\nabla f(0)| + |\nabla f|_{\text{Lip}}\,|x| \;\le\; \max\{|\nabla f(0)|,\,|\nabla f|_{\text{Lip}}\}\,\big(1 + |x|\big)\,, 
\end{equation} 
where $|\nabla f|_{\text{Lip}} := \sup_{x \neq y}\frac{\|\nabla f(x) - \nabla f(y)\|}{\|x - y\|}$ is the Lipschitz constant of $\nabla f$. Consequently (after some standard calculations on the SDE \eqref{eq:3.5}, see e.g. \cite{ks91}), we obtain that for all integers $p \ge 1$, 
\[ 
\mathbb{E}\!\Big[\,\max_{\,t \in [0,T]} \big|X_t^0\big|^{2p}\Big] < \infty\,.
\tag{3.8}
\] 

\noindent \textit{Proof of \eqref{eq:3.4}:} This limit is essentially established by stochastic calculus. Define $U_t := \exp\{\gamma^{-1}t\}\,V_t^\gamma$. Applying Itô’s formula and using the integration-by-parts formula (see Appendix), we have 
\[ 
dU_t = -\,\gamma^{-1}e^{\gamma^{-1}t}\,\nabla f(X_t^\gamma)\,dt + \sigma\,\gamma^{-1}e^{\gamma^{-1}t}\,dW_t\,,
\] 
or equivalently, 
\[ 
U_t = V_0 + \gamma^{-1}\int_{0}^{\,t} e^{\gamma^{-1}s}\,\nabla f(X_s^\gamma)\,ds + \sigma\,\gamma^{-1}\int_{0}^{\,t} e^{\gamma^{-1}s}\,dW_s\,.
\] 
It follows that 
\begin{equation}\label{eq:3.9}
\begin{aligned}
V_t^\gamma &= e^{-\gamma^{-1}t}\,U_t \\
&= e^{-\gamma^{-1}t}V_0 + \gamma^{-1}\int_{0}^{\,t} e^{\gamma^{-1}(s - t)}\,\nabla f(X_s^\gamma)\,ds + \sigma\,\gamma^{-1}\int_{0}^{\,t} e^{\gamma^{-1}(s - t)}\,dW_s\,,
\end{aligned}
\end{equation} 
and, further, 
\begin{align}
X_t^\gamma &= X_0^\gamma + \int_{0}^{\,t} V_s^\gamma\,ds \nonumber\\
&= X_0^\gamma + \int_{0}^{\,t} e^{-\gamma^{-1}s}ds\,V_0 \;-\; \gamma^{-1}\int_{0}^{\,t} \Bigg(\int_{s}^{\,t} e^{\gamma^{-1}(s-r)}dr\Bigg)\nabla f(X_s^\gamma)\,ds \nonumber\\ 
&\quad + \sigma\,\gamma^{-1}\int_{0}^{\,t} \Bigg(\int_{s}^{\,t} e^{\gamma^{-1}(s-r)}dr\Bigg) dW_s \nonumber\\ 
&= X_0^\gamma \;+\; \gamma\,V_0\,\big(1 - e^{-\gamma^{-1}t}\big) \;-\; \gamma^{-1}\int_{0}^{\,t} \Big(1 - e^{-\gamma^{-1}(t-s)}\Big)\nabla f(X_s^\gamma)\,ds \nonumber\\ 
&\quad + \sigma\,\gamma^{-1}\int_{0}^{\,t} \Big(1 - e^{-\gamma^{-1}(t-s)}\Big) dW_s\,,
\label{eq:3.10}
\end{align}
where in the second equality we used $\int_{0}^{\,t} e^{-\gamma^{-1}s}ds = \gamma\big(1 - e^{-\gamma^{-1}t}\big)$. Using Hölder’s inequality and the Lipschitz continuity of $\nabla f$, one can show that 
\[ 
\mathbb{E}\!\Bigg[\int_{0}^{\,t} 1_{\{X_s^\gamma = X_s^0\}}\,\nabla f(X_s^\gamma)\,ds\Bigg]^2 = 0\,.
\] 
Meanwhile, by \eqref{eq:3.9}, we have 
\[ 
\mathbb{E}\big[\gamma^2\,|V_0 - V_0^\gamma|^2\big] = o(\gamma^2)\,,
\] 
and hence, coming back to \eqref{eq:3.10} and applying Gronwall’s inequality, we obtain \eqref{eq:3.4}.

\medskip\noindent \textit{Proof of \eqref{eq:3.6}:} Before establishing \eqref{eq:3.6}, let us point out that, in contrast to some approaches (e.g. \cite{f04}), our proof does \emph{not} rely on a coupling argument. Instead, we leverage an exponential martingale and maximal inequalities. Define 
\[ 
I_t^\gamma := \int_{0}^{\,t} e^{-\gamma^{-1}(t-s)}\,dW_s\,,
\] 
so that $X_t^\gamma - X_t^0$ from \eqref{eq:3.10} can be written as 
\[ 
X_t^\gamma - X_t^0 = \gamma\,V_0\,\big(1 - e^{-\gamma^{-1}t}\big) - \gamma^{-1}\int_{0}^{\,t} \Big(1 - e^{-\gamma^{-1}(t-s)}\Big)\nabla f(X_s^\gamma)\,ds + \sigma\,\gamma^{-1} I_t^\gamma\,.
\] 
Taking the square and expectation and applying Jensen’s inequality, we deduce that 
\[ 
\mathbb{E}\!\Big[\,\max_{s\in[0,t]} |I_s^\gamma|\Big] \le C\sqrt{\gamma}\,,
\] 
for some constant $C$ depending only on $d$, $p$ and $\nabla f$. Moreover, by \eqref{eq:3.4}, 
\[ 
\int_{0}^{\,T_0} \mathbb{E}\!\big[\,|X_s^\gamma - X_s^0|\big]\,ds \le C\,\gamma\,,
\] 
and 
\[ 
\Bigg(\int_{0}^{\,T_0} \mathbb{E}\!\big[\,|X_s^\gamma - X_s^0|^2\big]\,ds\Bigg)^{\!1/2} \le C\,\gamma^{1/2}\,. 
\] 
Combining these estimates, we conclude that the rate $O(\gamma^{1/2})$ leading to \eqref{eq:3.6} indeed holds. 
\end{proof}

\section{Appendix}

\noindent \textbf{Gronwall’s lemma:} Let $f$, $\alpha$, $g$ be three real-valued functions defined on an interval $[0,T]$ such that 
\[ 
\forall\,t \in [0,T]: \qquad f(t) \le g(t) + \int_{0}^{\,t} \alpha(s)\,f(s)\,ds\,,
\] 
and assume $g$ is non-decreasing and $\alpha(s) \ge 0$. Then 
\[ 
\forall\,t \in [0,T]: \qquad f(t) \le g(t)\,\exp\!\Big\{\int_{0}^{\,t} \alpha(s)\,ds\Big\}\,. 
\]

\noindent \textbf{Covariation (cross-variation) and quadratic variation:} If $X$ and $Y$ are two real-valued continuous semimartingales given by 
\[ 
X_t = X_0 + M^X_t + A^X_t, \qquad Y_t = Y_0 + M^Y_t + A^Y_t, \qquad t \ge 0, 
\] 
with $M^X$ and $M^Y$ the respective martingale parts and $A^X$ and $A^Y$ the finite-variation parts, then the covariation (cross-variation) process of $X$ and $Y$ is the process $\langle X, Y \rangle = (\langle X,Y\rangle_t)_{t\ge0}$ defined by 
\[ 
\langle X, Y \rangle_t := \lim_{\max |t_{i+1}-t_i|\to\,0}\;\sum_{k=0}^{N-1} \big(X_{t_{k+1}} - X_{t_k}\big)\,\big(Y_{t_{k+1}} - Y_{t_k}\big)\,. 
\] 
In particular, $\langle X \rangle := \langle X, X\rangle$ corresponds to the quadratic variation process of $X$.

\noindent \textbf{Itô’s integration by parts:} Let $(X_t)_{t\ge0}$ and $(Y_t)_{t\ge0}$ be two continuous semimartingales. Then, almost surely for all $t \ge 0$, 
\[ 
X_t Y_t = X_0 Y_0 + \int_{0}^{\,t} X_s\,dY_s + \int_{0}^{\,t} Y_s\,dX_s + \langle X, Y \rangle_t\,. 
\] 
In particular, if $X_t = \int_{0}^{\,t} \sigma_s\,dW_s$ and $Y_t = \int_{0}^{\,t} b_s\,ds$, where $\sigma$ and $b$ are integrable $\,\mathcal{F}_t$-adapted processes, the above gives 
\[ 
X_t Y_t = \int_{0}^{\,t} X_s\,b_s\,ds + \int_{0}^{\,t} Y_s\,\sigma_s\,dW_s, \qquad t \ge 0\,.
\] 

\noindent \textbf{Itô’s formula for general semimartingales:} Let $(X_t)_{t\ge0}$ be an $\mathbb{R}^d$-valued continuous semimartingale with the canonical decomposition $X_t = X_0 + M_t + A_t$ (where $M$ is a martingale and $A$ is a finite-variation process). Then, for any scalar function $f: [0,T]\times\mathbb{R}^d \to \mathbb{R}$ satisfying the regularity conditions: 
\begin{itemize}
\item for each fixed $x$, the map $t \mapsto f(t,x)$ is continuously differentiable on $[0,T]$ (with continuous $\partial_t f$); 
\item for each fixed $t$, the map $x \mapsto f(t,x)$ is $C^2$ on $\mathbb{R}^d$ (with continuous second partial derivatives), 
\end{itemize}
it holds that 
\[
\begin{aligned}
f(t,X_t) &= f(0,X_0) + \int_{0}^{\,t} \partial_s f(s, X_s)\,ds \\
&\quad + \sum_{i=1}^d \int_{0}^{\,t} \partial_{x_i} f(s,X_s)\,dX_s^i \\
&\quad + \frac{1}{2}\sum_{i,j=1}^d \int_{0}^{\,t} \partial^2_{x_i x_j} f(s,X_s)\,d\langle X^i, X^j \rangle_s\,.
\end{aligned}
\] 
\noindent \textbf{Equivalent formulation:} Writing $\nabla_x f$ and $\nabla_x^2 f$ for the gradient and Hessian of $f$, and $\langle X, X \rangle_t$ for the $d\times d$ covariance matrix of $X$ (i.e. $\langle X,X\rangle_t^{(i,j)} = \langle X^i, X^j\rangle_t$), we can express the Itô formula as 
\[
\begin{aligned}
f(t,X_t) &= f(0,X_0) + \int_{0}^{\,t} \partial_s f(s,X_s)\,ds \\
&\quad + \int_{0}^{\,t} \nabla_x f(s,X_s)\cdot dM_s + \int_{0}^{\,t} \nabla_x f(s,X_s)\cdot dA_s \\
&\quad + \frac{1}{2} \int_{0}^{\,t} \text{Tr}\!\Big[\nabla_x^2 f(s,X_s)\,d\langle X,X\rangle_s\Big]\,.
\end{aligned}
\] 
(This formulation shows, under general assumptions, that the process $f(t, X_t)$ of a semimartingale $X$ remains a semimartingale.)

\vfill

\printbibliography

% Ensure all references are included
\nocite{Hottovy2015}

\end{document}
